\begin{lemma}
For every $\eta>0$ there are $\Delta_0$ and $\gamma_0$ such that the
following holds.  Let $G$ be a $\Delta$-regular graph with
$\Delta\ge\Delta_0$, let $\gamma\ge\gamma_0$, and let $I$ be sampled by
\noderef{RandomIndependentSetSampling}.  For every vertex $r$ and every
$X\subseteq N_G(r)$ such that no two distinct vertices of $X$ have a common
neighbor outside $N_G[r]$,
\[
\Pr[I\cap X\ne\emptyset]\le b_{G[X\cup\{r\}],\Delta}(r)+\eta ,
\]
where the local parameter of the induced graph is the quantity defined in
\noderef{SamplingInducedNeighborhoodLocalBParameter}.
\end{lemma}
\begin{proof}
Let $I_X=I\cap X$, let
$\overline P_X=\binom{|I_X|}{2}$, and let
$\overline T_X=\binom{|I_X|}{3}$.  The third Bonferroni inequality applied
to the events $\{x\in I\}$ for $x\in X$ gives
\[
\Pr[I_X\ne\emptyset]
\le \mathbb E|I_X|-\mathbb E\overline P_X+\mathbb E\overline T_X .
\]
Let $\mathcal J_2$ and $\mathcal J_3$ be the independent pairs and triples in
$X$.  Since every vertex of $X$ is adjacent to $r$, these are exactly the
two- and three-element independent subsets of the neighborhood of $r$ in the
induced graph $G[X\cup\{r\}]$.  Therefore
$|\mathcal J_2|=P_{G[X\cup\{r\}]}(r)$ and
$|\mathcal J_3|=T_{G[X\cup\{r\}]}(r)$ by
\noderef{IndependentPairCount} and \noderef{IndependentTripleCount}.

Choose $\Delta_0$ and $\gamma_0$ large enough for
\noderef{SamplingPairTripleBonferroniCore} and
\noderef{SamplingOneVertexEstimate} with the error budgets used below.
Applying \noderef{SamplingPairTripleBonferroniCore} with error budget
$\eta/2$ gives
\[
\mathbb E\overline P_X-\mathbb E\overline T_X
\ge
\frac{|\mathcal J_2|}{\Delta^2}
-\frac{|\mathcal J_3|}{\Delta^3}
-\eta/2 .
\]

It remains to bound the first moment.  If $X=\emptyset$, then
$I_X=\emptyset$, $\mathcal J_2=\mathcal J_3=\emptyset$, and
$b_{G[X\cup\{r\}],\Delta}(r)=0$; the desired probability is then $0$.
Assume $X$ is nonempty.  For each $x\in X$, the identity in
\noderef{SamplingOneVertexEstimate} gives
\[
\Pr[x\in I]
=\frac1\Delta\int_0^\gamma \left(1-\frac t\Delta\right)^\Delta\,dt .
\]
As in the proof of that node, the sampling law gives $0<\gamma\le\Delta$ in
the non-vacuous case.  Hence
$0\le(1-t/\Delta)^\Delta\le e^{-t}$ for $0\le t\le\gamma$, and
\[
\Pr[x\in I]\le
\frac1\Delta\int_0^\gamma e^{-t}\,dt
\le\frac1\Delta .
\]
Summing over $X$ yields
$\mathbb E|I_X|=\sum_{x\in X}\Pr[x\in I]\le |X|/\Delta$.

Combining this first-moment bound with Bonferroni and the pair-minus-triple
core estimate gives
\[
\Pr[I_X\ne\emptyset]\le
\frac{|X|}{\Delta}-\frac{|\mathcal J_2|}{\Delta^2}
+\frac{|\mathcal J_3|}{\Delta^3}+\eta/2 .
\]
By \noderef{LocalBParameter} and the identifications supplied by
\noderef{IndependentPairCount} and \noderef{IndependentTripleCount}, the
right-hand side is
$b_{G[X\cup\{r\}],\Delta}(r)+\eta/2$, which is at most
$b_{G[X\cup\{r\}],\Delta}(r)+\eta$ because $\eta>0$.  This proves the
claim.
\end{proof}
